\section*{अध्याय \ref{ch:b3:quantum-model-systems} --- रसायनज्ञों के लिए क्वांटम यांत्रिकी: मॉडल निकाय}

\begin{solution}{exo:b3:quantum-model-systems:1}
$E_1 = h^2/8m_eL^2$। (a) $L = \qty{1.00}{nm}$: $E_1 = \qty{6.02e-20}{J} =
\qty{0.376}{eV}$; $\Delta E_{12} = 3E_1$ और $\lambda = hc/3E_1 =
\qty{1099}{nm}$, निकट अवरक्त में। (b) $L = \qty{0.50}{nm}$: $E_1$ चार
गुना बड़ा है, \qty{2.41e-19}{J} $= \qty{1.50}{eV}$, और $\lambda =
\qty{275}{nm}$, पराबैंगनी में।
\end{solution}

\begin{solution}{exo:b3:quantum-model-systems:2}
$n_x^2 + n_y^2 + n_z^2 = 3$ (1,1,1): $g = 1$; $6$ (2,1,1 और उसके क्रमचय):
$g = 3$; $9$ (2,2,1): $g = 3$; $11$ (3,1,1): $g = 3$; $12$ (2,2,2): $g = 1$।
\end{solution}

\begin{solution}{exo:b3:quantum-model-systems:3}
$[\hat x^2,\hat p]f = -\iu\hbar x^2f' + \iu\hbar(x^2f)' = 2\iu\hbar xf$, अतः
$[\hat x^2,\hat p] = 2\iu\hbar\hat x$। $[\hat p^2,\hat x]f = -\hbar^2[(xf)'' -
xf''] = -2\hbar^2f'$, अतः $[\hat p^2,\hat x] = -2\iu\hbar\hat p$।
\end{solution}

\begin{solution}{exo:b3:quantum-model-systems:4}
$E_J/hc = BJ(J+1)$: 0, 21.19, 63.56 और \qty{127.12}{cm^{-1}}, जहाँ $g = 1$,
3, 5, 7। $I = h/(8\pi^2cB) = \qty{2.643e-47}{kg.m^2}$ ($c$ को
\unit{cm/s} में लेकर)।
\end{solution}

\begin{solution}{exo:b3:quantum-model-systems:5}
$|\psi_n|^2$, $L/2$ के सापेक्ष सममित है, अतः $\langle x\rangle = L/2$।
$\sin^2u = \frac12(1 - \cos2u)$ और दो बार खंडशः समाकलन से
$\langle x^2\rangle = L^2\big(\frac13 - \frac{1}{2n^2\pi^2}\big)$। $n = 1$ के लिए
$\langle x^2\rangle = 0.283L^2$, जो चिरसम्मत $L^2/3$ से कम है: मूल
अवस्था बीच में जमा होती है। $n$ बढ़ने पर चिरसम्मत मान पुनः प्राप्त होता है।
\end{solution}

\begin{solution}{exo:b3:quantum-model-systems:6}
$1\,\unit{cm^{-1}} = \qty{0.011963}{kJ/mol}$। \omterm{def:b3:quantum-model-systems:oscillator}{शून्य-बिंदु ऊर्जा}(\ce{H2}) $= 2200.6 \times
0.011963 = \qty{26.33}{kJ/mol}$, \omterm{def:b3:quantum-model-systems:oscillator}{शून्य-बिंदु ऊर्जा}(\ce{D2}) $= \qty{18.63}{kJ/mol}$; अंतर
\qty{7.69}{kJ/mol}। \omterm{def:b3:quantum-model-systems:oscillator}{बल स्थिरांक} समान है, अतः $\tilde\omega \propto \mu^{-1/2}$
और अपेक्षित अनुपात $\sqrt{\mu_D/\mu_H} = 1.4137$ है (परमाणु द्रव्यमान
2.01410 और 1.00783); मापा गया अनुपात 1.4127 है, 0.07\,\% के भीतर (यह
छोटा अंतर इलेक्ट्रॉनों से आता है, जो नाभिकों का पूर्णतः
अनुसरण नहीं करते)।
\end{solution}

\begin{solution}{exo:b3:quantum-model-systems:7}
$\int_0^Lx^2(L - x)^2\,\dd x = L^5/30$, अतः $N = \sqrt{30/L^5}$। तब
\[
\langle\psi_1|\phi\rangle = \sqrt{\tfrac2L}\sqrt{\tfrac{30}{L^5}}\int_0^Lx(L -
x)\sin\tfrac{\pi x}{L}\,\dd x = \frac{\sqrt{60}}{L^3}\cdot\frac{4L^3}{\pi^3} =
\frac{4\sqrt{60}}{\pi^3} = 0.99928 .
\]
परवलय 99.86\,\% मूल अवस्था है (अतिव्यापन का वर्ग)।
\end{solution}

\begin{solution}{exo:b3:quantum-model-systems:8}
$\langle f|\hat pg\rangle = -\iu\hbar\int f^*g'\,\dd x = -\iu\hbar[f^*g] +
\iu\hbar\int f^{*\prime}g\,\dd x = \int(-\iu\hbar f')^*g\,\dd x = \langle\hat
pf|g\rangle$, क्योंकि कोष्ठक सिरों पर लुप्त होता है। $\hat x\hat p$ के लिए, इस तथ्य का उपयोग करते हुए कि $\hat x$ और $\hat p$ दोनों \omterm{def:b3:quantum-model-systems:hermitian}{हर्मिटी} हैं:
$\langle f|\hat x\hat pg\rangle = \langle\hat xf|\hat pg\rangle = \langle\hat p\hat
xf|g\rangle$, और $\hat p\hat x = \hat x\hat p - \iu\hbar \ne \hat x\hat p$: यह
\omterm{def:b3:quantum-model-systems:hermitian}{हर्मिटी} नहीं है (उसका सममितीकृत रूप $\frac12(\hat x\hat p + \hat p\hat x)$ \omterm{def:b3:quantum-model-systems:hermitian}{हर्मिटी} है)।
\end{solution}

\begin{solution}{exo:b3:quantum-model-systems:9}
$L = 6 \times 139.7 = \qty{838}{pm}$, $N = 6$: $\lambda = 8m_ecL^2/(7h) =
\qty{331}{nm}$। अवशोषण पराबैंगनी में है: हेक्साट्राईन
रंगहीन है।
\end{solution}

\begin{solution}{exo:b3:quantum-model-systems:10}
$\kappa = \sqrt{2m(V_0 - E)}/\hbar$, जहाँ $V_0 - E = \qty{0.20}{eV}$:
$\kappa_H = \qty{9.82e10}{m^{-1}}$, $\kappa_D = \qty{1.388e11}{m^{-1}}$। अनुपात में
पूर्वगुणक निरस्त हो जाते हैं: $T_H/T_D = \eu^{2a(\kappa_D - \kappa_H)} =
\eu^{4.06} = 58$, यथार्थ मान: अवरोध दोनों के लिए मोटा है
($\kappa a = 4.9$ और 6.9)।
\end{solution}

\begin{solution}{exo:b3:quantum-model-systems:11}
स्तर $E_m = m^2\hbar^2/2m_eR^2$: $m = 0$ में दो इलेक्ट्रॉन आते हैं, $m = \pm1$ में
चार। उच्चतम भरा स्तर $|m| = 1$ है, निम्नतम रिक्त $|m| = 2$: $\Delta E =
3\hbar^2/2m_eR^2 = \qty{9.38e-19}{J}$, $\lambda = \qty{212}{nm}$,
पराबैंगनी में, जहाँ बेंज़ीन का प्रबल अवशोषण होता भी है।
\end{solution}

\begin{solution}{exo:b3:quantum-model-systems:12}
$\langle r\rangle = \frac{4\pi}{\pi a^3}\int_0^\infty r^3\eu^{-2r/a}\dd r =
\frac{4}{a^3}\cdot\frac{3!}{(2/a)^4} = \frac32a$, अर्थात् \qty{79.4}{pm}, जब $a = a_0$।
$\langle 1/r\rangle = \frac{4}{a^3}\cdot\frac{1!}{(2/a)^2} = \frac1a$, अतः
$\langle V\rangle = -e^2/(4\pi\varepsilon_0a)$। चूँकि $E_1 =
-e^2/(8\pi\varepsilon_0a)$ ($a = 4\pi\varepsilon_0\hbar^2/\mu e^2$ से),
$\langle V\rangle = 2E_1$, और माध्य गतिज ऊर्जा $-E_1$ है।
\end{solution}

\begin{solution}{pb:b3:quantum-model-systems:1}
\textbf{1.} $j$ परमाणुओं में से हर एक एक p कक्षक देता है; $j - 1$ कार्बन और
एक नाइट्रोजन एक-एक इलेक्ट्रॉन देते हैं और दूसरा नाइट्रोजन अपना एकाकी युग्म,
क्योंकि धनायन का आवेश शृंखला पर फैला है: $N = j + 1$, अर्थात् 6,
8 और 10।
\textbf{2.} छह इलेक्ट्रॉन $n = 1$, 2, 3 को भरते हैं: HOMO $n = 3$, LUMO $n = 4$।
\textbf{3.} $\Delta E = E_{N/2+1} - E_{N/2} = (N+1)h^2/8mL^2$ और $\lambda =
hc/\Delta E = 8mcL^2/h(N+1)$।
\textbf{4.} $L_0 = 4l$, $6l$, $8l$: 559, 838 और \qty{1118}{pm}।
\textbf{5.} 147, 257 और \qty{374}{nm}।
\textbf{6.} सभी बहुत छोटे (3.6, 2.3 और 1.9 के गुणकों से): मॉडल का डिब्बा बहुत
छोटा है, क्योंकि $\lambda \propto L^2$; इलेक्ट्रॉन नाइट्रोजनों के बीच की शृंखला से
लंबी दूरी पर गति करते हैं।
\textbf{7.} मूलबिंदु केंद्र पर रखने पर, $\psi_n \propto \cos(n\pi u/L)$
विषम $n$ के लिए और $\sin(n\pi u/L)$ सम $n$ के लिए ($u = x - L/2$): ये
$u$ के सम और विषम फलन हैं।
\textbf{8.} यदि $n + n'$ सम है, तो $\psi_n$ और $\psi_{n'}$ की समता समान है;
उनका गुणनफल गुणा $u$ विषम है और सममित अंतराल पर उसका
समाकल शून्य है।
\textbf{9.} HOMO $N/2$ और LUMO $N/2 + 1$ के लिए $n + n' = N + 1$, विषम: सदा
अनुमत।
\textbf{10.} दूसरा रंजक: HOMO 4, LUMO 5। $4 \to 6$: योग सम, वर्जित।
$3 \to 5$: योग सम, वर्जित।
\textbf{11.} $n = 4$ से अगला अनुमत संक्रमण $4 \to 7$ है, जिसका
$\Delta E$ $(49 - 16)/(25 - 16) = 3.67$ गुना बड़ा है: $\lambda/3.67$ पर, अर्थात्
दूसरे रंजक के लिए \qty{164}{nm}, सुदूर पराबैंगनी में।
\textbf{12.} यह केवल शृंखला के केंद्र के सापेक्ष विभव की सममिति पर निर्भर है,
जो वास्तविक सममित रंजकों में भी है।
\textbf{13.} $E = hc/\lambda = \qty{3.294e-19}{J} = \qty{2.06}{eV}$।
\textbf{14.} $hc\tilde\omega_e = \qty{0.371}{eV}$।
\textbf{15.} $2hcB = \qty{4.21e-22}{J} = \qty{2.63}{meV}$।
\textbf{16.} $kT = \qty{25.7}{meV}$।
\textbf{17.} केवल घूर्णी स्तर ($2hcB \ll kT$); कंपनिक ($15kT$) और
इलेक्ट्रॉनिक ($80kT$) उत्तेजन व्यावहारिक रूप से समष्टित नहीं होते।
\textbf{18.} इलेक्ट्रॉनिक: दृश्य और पराबैंगनी; कंपनिक: अवरक्त;
घूर्णी: सूक्ष्मतरंगें (और सुदूर अवरक्त)।
\textbf{19.} $L = \sqrt{\lambda h(N+1)/8mc}$, जहाँ $N = 8$: $L = \qty{1283}{pm}$।
\textbf{20.} $\delta = (1283 - 838)/2 = \qty{222}{pm}$।
\textbf{21.} $L = 559 + 445 = \qty{1003}{pm}$, $\lambda = \qty{474}{nm}$, 524 nm से 9.5\,\%
नीचे।
\textbf{22.} $L = 1118 + 445 = \qty{1562}{pm}$, $\lambda = \qty{732}{nm}$, 711 nm से 2.9\,\%
ऊपर।
\textbf{23.} उन ऐरोमैटिक वलयों में जिन पर नाइट्रोजन हैं: $\pi$
तंत्र नाइट्रोजन परमाणुओं पर समाप्त नहीं होता, और प्रभावी डिब्बे में
हर वलय का कुछ भाग सम्मिलित है।
\textbf{24.} $\delta/l = 222/139.7 = 1.6$: \textbf{डिब्बा हर नाइट्रोजन से आगे
लगभग \qty{222}{pm}, अर्थात् डेढ़ आबंध, तक फैला है}, और उस एक
लंबाई से मॉडल तीनों रंगों को 10\,\% के भीतर पुनः उत्पन्न करता है।
\end{solution}
